\newpage
\titledquestion{Balanced?}

\textbf{Balanced?}

A tree is \textbf{balanced} if at every node its two subtrees' heights differ by
at most one, and both subtrees are balanced.
\begin{codeblock}
  datatype tree = Empty | Node of tree * int * tree
\end{codeblock}

Here is the natural first implementation, using a helper \code{height}:
\begin{codeblock}
  fun height Empty = 0
    | height (Node (l, _, r)) = 1 + Int.max (height l, height r)

  fun isBalanced Empty = true
    | isBalanced (Node (l, _, r)) =
        abs (height l - height r) <= 1
        andalso isBalanced l
        andalso isBalanced r
\end{codeblock}

Throughout, assume the balance check passes (so the function fully recurses),
the tree is balanced with \code{n} nodes, and \code{height T} has work $O(m)$
for \code{m} nodes.

\begin{parts}
  \part[4]\hfill

    Write and solve the \textbf{work} recurrence for the naive \code{isBalanced}.

    \begin{solutionorbox}[14em]
      At each node the two \code{height} calls walk the whole subtree, costing
      $O(n)$, and we recur on both halves:
      $$W(n) = 2W(n/2) + O(n).$$
      By the recursion tree each of the $\log n$ levels does $O(n)$ work, so
      $$W(n) = O(n \log n).$$
    \end{solutionorbox}

  \part[4]\hfill

    Now take a \textbf{left spine}: every node has only a left child. (Pretend the
    balance check still passes, so the function recurses all the way down.) Write
    and solve the \textbf{work} recurrence for the naive \code{isBalanced} on a
    spine of \code{n} nodes, and say in one sentence why it is so much worse.

    \begin{solutionorbox}[14em]
      There is one recursive call per node, and \code{height} at the top walks all
      \code{n} nodes, then $n-1$, and so on:
      $$W(n) = W(n-1) + O(n) = O(n) + O(n-1) + \cdots = O(n^2).$$

      On a spine \code{height} rescans the entire remaining chain at every level,
      so the naive version is \textbf{quadratic} --- not merely a $\log n$ factor
      worse than the balanced case.
    \end{solutionorbox}

  \newpage

  \part[8]\hfill

    Implement \code{isBalanced} so that it visits each node only once, computing
    each subtree's height and balancedness in a single pass.

    \textbf{Hint:} a helper \code{check : tree -> bool * int} returning both,
    so no height is recomputed.

    \begin{solutionorbox}[20em]
      \begin{codeblock}
        fun check Empty = (true, 0)
          | check (Node (l, _, r)) =
              let
                val (bl, hl) = check l
                val (br, hr) = check r
              in
                ( bl andalso br
                    andalso abs (hl - hr) <= 1,
                  1 + Int.max (hl, hr) )
              end

        fun isBalanced T = #1 (check T)
      \end{codeblock}

      Each node is visited once and does $O(1)$ work: it recurses into each child
      \textit{once}, getting back both the child's balancedness and its height,
      and combines them. No separate \code{height} traversal is ever needed.
    \end{solutionorbox}

  \part[4]\hfill

    Write and solve the \textbf{work} recurrence for your single-pass version, and
    compare it to the naive version on both shapes.

    \begin{solutionorbox}[12em]
      Each node does $O(1)$ work and we recur on both halves:
      $$W(n) = 2W(n/2) + O(1) = O(n).$$
      Every node is visited exactly once --- so this holds on \textbf{any} shape,
      including the spine, beating the naive $O(n \log n)$ and $O(n^2)$.
    \end{solutionorbox}

  \part[4]\hfill

    Write and solve the \textbf{span} recurrence for your single-pass version.
    Why does it differ from the work?

    \begin{solutionorbox}[13em]
      $$S(n) = S(n/2) + O(1) = O(\log n).$$

      The $2$ becomes a $1$: the recursive calls \code{check l} and \code{check r}
      are \textbf{independent}, so they run \textbf{in parallel} and the span
      follows only the deeper subtree, rather than summing both as the work does.
    \end{solutionorbox}
\end{parts}
